\begin{table}[t]
\tbl{\label{tab:spreads:define}}
{\begin{tcolorbox}[tabularx={p{3cm}|p{2cm}|p{6.3cm}|},enhanced,fonttitle=\bfseries\large,fontupper=\normalsize\sffamily,
colback=yellow!10!white,colframe=red!50!black,colbacktitle=blue!30!white,
coltitle=black,title=Creating a Spread]
\hline
Command & Arguments & Purpose \\
\hline
\hline
\verb'-kernel_field'  &  $F$  &  Define the kernel of the spread. $F$ must be an object of type finite field.\\
\hline
\verb'-group'  &  $G$  &  Define the group acting on the spread. Should be $\PGGL(2k,F)$. \\
\hline
\verb'-group_on_subspaces'  &  $G$  &   \\
\hline
\verb'-k'  &  $k$  &  Set the dimension of the spread. \\
\hline
\verb'-catalogue'  &  $i$  &  Pull spread number $i$ from the catalogue of spreads associated with the given field and the given dimension. \\
\hline
\verb'-family'  &  $L$  & Define a spread from a named family $L.$ So far, no family has been provided. \\
\hline
\verb'-spread_set'  &  label-txt label-tex $S$  & Define a spread from the named spreadset $S$. The spreadset $S$ must be a vector object. It must contain $q^3k^2$ entries over $F.$ \\
\hline
\verb'-transform'  &  elt  &  Apply the transformation given by the group element. \\
\hline
\verb'-transform_inverse'  &  elt  &  Apply the inverse transformation given by the group element. \\
\hline
\end{tcolorbox}}
\end{table}
